\begin{frame}{핵심 숫자 ①·②: 표본 잡음과 학습 곡선}
  \small
  \begin{columns}[T]
    \begin{column}{0.5\textwidth}
      \kb{① 정확도 차이는 \textbf{표본 잡음보다 작다}}
      \begin{itemize}
        \item 세 모델 test: logreg \textbf{0.959}, MLP
          \textbf{0.953}, GBDT \textbf{0.942} —
          \textbf{max$-$min 0.018}
        \item \textbf{같은 모델·가중치}로 test 100건을 60번
          바꾸면 \textbf{0.930$\sim$0.990}(범위 \textbf{0.06})
        \item 모델 간 차이(0.018)는 표본 흔들림(0.06)의
          \textbf{1/3} — ``MLP 1위''는 \textbf{표본 선택}에
          의존
      \end{itemize}
    \end{column}
    \begin{column}{0.5\textwidth}
      \kb{② 학습 곡선이 \textbf{에폭 수}를 정한다}
      \begin{itemize}
        \item val은 에폭 \textbf{25$\sim$32 포화}, 그 뒤에도
          \textbf{train 손실만} 하락 $\to$ 과적합 신호
        \item ``에폭 50''이 \textbf{포화(≈30)} 근거라면 3점
          통과, \textit{시간 제약} 숫자라면 ``30에서 멈췄을
          때 val이 다른가'' 증명 — \textbf{8.1 M3의 이유}
      \end{itemize}
    \end{column}
  \end{columns}
\end{frame}
